\begin{RequAnswer}{1.1}
A proposition is a statement that is either true or false.
\end{RequAnswer}
\begin{RequAnswer}{1.2}
The operators are {\em negation}, {\em or}, {\em and}, {\em exclusive-or}, {\em conditional}, and {\em biconditional}.
See Table~\ref{tab:not} and  Table~\ref{tab:2vars} for the truth tables.
\end{RequAnswer}
\begin{RequAnswer}{1.3}
$p\vee q$ is true if $p$ is true, $q$ is true, or both $p$ and $q$ are true, whereas
$p\oplus q$ is true if and only if exactly one of $p$ or $q$ is true. So the difference is
that is both $p$ and $q$ are true, $p\vee q$ is true, but $p\oplus a$ is false.
\end{RequAnswer}
\begin{RequAnswer}{1.4}
It is true unless $p$ is true and $q$ is false.
\end{RequAnswer}
\begin{RequAnswer}{1.5}
Here is a truth table with some intermediate columns to help. As long as you have the same final column, yours is probably fine.

$\begin{array}{|cc||c|c|c|} \hline p & q & p  \wedge q & \neg p&(p \wedge q) \vee\neg p \\
\hline
T & T & T & F&T\\
T & F & F & F&F \\
F & T & F & T&T \\
F & F & F & T&T\\
\hline
\end{array}$
\end{RequAnswer}
\begin{RequAnswer}{1.6}
Here is a truth table with some intermediate columns to help. As long as you have the same final column, yours is probably fine.

$\begin{array}{|lll|l|l|}
\hline
p & q & r & p\wedge q& (p\wedge q)\vee r\\
\hline
T & T & T & T &T\\
T & T & F & T &T\\
T & F & T & F &T\\
T & F & F & F &F\\
F & T & T & F &T\\
F & T & F & F &F\\
F & F & T & F &T\\
F & F & F & F &F\\
\hline
\end{array}$
\end{RequAnswer}
\begin{RequAnswer}{1.7}
By definition, no. If $p$ is true, $\neg p$ is false, and if $p$ is false, $\neg p$ is true.
\end{RequAnswer}
\begin{RequAnswer}{1.8}
$q$ must be true since one of them has to be and $p$ is not.
\end{RequAnswer}
\begin{RequAnswer}{1.9}
Nothing. It might be true, but it also might be false.
\end{RequAnswer}
\begin{RequAnswer}{1.10}
They are both true.
\end{RequAnswer}
\begin{RequAnswer}{1.11}
$q$ has to also be false since $p\leftrightarrow q$ implies they have the same truth value.
\end{RequAnswer}
\begin{RequAnswer}{1.12}
In this case $q$ also has to be true.
\end{RequAnswer}
\begin{RequAnswer}{1.13}
Since $p\rightarrow q$ is true whenever $p$ is false, we cannot say anything about $q$,
so it could be true or false.
\end{RequAnswer}
\begin{RequAnswer}{1.14}
By definition, if a proposition is a contingency, then it is {\em not}
a tautology.
\end{RequAnswer}
\begin{RequAnswer}{1.15}
By definition, a proposition and its negation can {\em never} both be true.
If the proposition is true, its negation is false, and if the proposition is false, its negation is true.
\end{RequAnswer}
\begin{RequAnswer}{1.16}
(a) $p\vee T$ is true if and only if either $p$ or $T$ is true. Since clearly $T$ is true, $p\vee T$ is always true.
Therefore, $ p\vee T \ = T$.
Alternatively, you can give a truth table for $p\vee T$ and then comment something like
``since the final column of the truth table is always true, $ p\vee T \ = T$.''
(b) We will do this one with a truth table:
Notice that the column for $p\wedge F$ in the truth table below is always false. Therefore  $p\wedge F = F$.

$\begin{array}{|c||c|} \hline p & p\wedge F \\
\hline
T & F\\
F & F\\
\hline
\end{array}$
\end{RequAnswer}
\begin{RequAnswer}{1.17}
You could draw a truth table and show that the columns for these two differ. However, there is any easier approach.
Notice that if $p$ is true and $q$ is false, $\neg p \wedge \neg q$ is false, but $\neg(p\wedge q)$ is true.
Therefore they are not equivalent.
\end{RequAnswer}
\begin{RequAnswer}{1.18}
We can't use the first reason because DeMorgan's law isn't the
only rule we have to determine whether or not two propositions
are equivalent. Perhaps they are equivalent by some other rule
(in this case they aren't, which you should know if you
answered the previous question).
The second reason is also not valid because
every proposition is equivalent to many other
propositions that look different than it, so just knowing that
they look different doesn't tell us anything. For instance,
DeMorgan's law says that $\neg p \wedge \neg q=\neg(p\vee q)$.
So even though those two don't look that same, they are equivalent.
\end{RequAnswer}
\begin{RequAnswer}{1.19}
Find an assignment of true values to the variables
such that one is true and the other is false. That's all there is to it.
Although sometimes this is easier said than done!
\end{RequAnswer}
\begin{RequAnswer}{1.20}
A propositional function is a statement that contains one or more variables and depending on the values of the variable(s),
the statement is true or false. In other words, a propositional function is a function whose outputs are propositions.
\end{RequAnswer}
\begin{RequAnswer}{1.21}
$\neg\forall x P(x)$ means that it is not the case that $P(x)$ is true for all values of $x$.
So it doesn't mean it is never true. It means that for one or more values of $x$ it is false.
That is, it means that it is not always true.
\end{RequAnswer}
\begin{RequAnswer}{1.22}
$\neg\exists x P(x)$ means that it is not the case that there is a value of $x$ for which $P(x)$ is true.
In other words, it is indeed saying that $P(x)$ is never true.
\end{RequAnswer}
\begin{RequAnswer}{1.23}
The two obvious alternatives are
$\neg\exists x  \exists y Q(x,y)$ and
$\forall x \forall y\neg  Q(x,y)$.
\end{RequAnswer}
\begin{RequAnswer}{1.24}
These are all equivalent: $ \forall x  \neg(x<0 \wedge x>0)$,
$ \forall x  (\neg(x<0) \vee  \neg(x>0))$, and
$ \forall x  (x>=0 \vee  x<=0)$.
\end{RequAnswer}
\begin{RequAnswer}{1.25}
$\forall x\exists y H(x,y)$.
\end{RequAnswer}
\begin{RequAnswer}{1.26}
Let $C(x,y)=$``$x$ changes at time $y$.''
Then $\neg C(x,y)=$``$x$ stays the same at time $y$.''
{\em If you did not define a predicate similar to this, stop reading now
and try to come up with the rest of the answer using $C(x,y)$ before continuing to read.}
Then the expression can be written as $\neg\exists x\exists y  C(x,y)\wedge \neg\exists x\exists y \neg C(x,y)$,
which of course can be simplified $\forall x\forall y \neg C(x,y) \wedge \forall x \forall y C(x,y)$
(but you would probably re-interpret this version in English as ``Everything stays the same, everything changes,''
which if you think about is it essentially saying the same thing.).
\end{RequAnswer}
\begin{RequAnswer}{1.27}
(a) Given any integer, there is an integer that is at least as large as it.
(b) There is an integer such that every integer is at least as large as it is.
(c) They do not seem to be saying the same thing at all.
(d) True.
(e) False.
(f) If the universe of discourse is changed to positive integers, then both statements are true.
(The second one becomes true because every positive integer is at least as large as 1.)
(g) No. Just because two statements have the same truth value, that does not mean they are saying the same thing.
For instance, ``1+1=2" is true and ``$x^2\geq 0$'' is true, but they definitely are saying very different things.
\end{RequAnswer}
\begin{RequAnswer}{1.28}
$\neg p$, \, $q$, \, $\neg r$, \, $r$.
\end{RequAnswer}
\begin{RequAnswer}{1.29}
$\neg p$, \, $p\wedge r$, \, $\neg p \wedge r$, \, $q$,  \, $\neg r$.
\end{RequAnswer}
\begin{RequAnswer}{1.30}
$\neg p$, \,
$p\wedge r$, \,
$q\vee \neg r$, \,
$\neg p \wedge r$, \,
$(p\wedge q)\vee (q\wedge \neg r)\vee \neg p$, \,
$(p\wedge \neg r\wedge q)\vee (\neg p\wedge r \wedge \neg q)\vee(p\wedge r \wedge q)\vee(\neg p\wedge \neg r \wedge \neg q)$.
\end{RequAnswer}
